\documentclass[11pt,a4paper,oneside]{report}

% DevSecOps process report for the Timesheet application.
% Compile twice with pdfLaTeX. Personal screenshots are optional and have
% explicit insertion points under images/own/.
\usepackage[T1]{fontenc}
\usepackage[utf8]{inputenc}
\usepackage[french]{babel}
\usepackage{lmodern}
\usepackage{microtype}
\usepackage[a4paper,margin=2.35cm,headheight=15pt]{geometry}
\usepackage{graphicx}
\usepackage{float}
\usepackage{booktabs}
\usepackage{tabularx}
\usepackage{longtable}
\usepackage{array}
\usepackage{enumitem}
\usepackage{xcolor}
\usepackage{tikz}
\usetikzlibrary{arrows.meta,positioning,shapes.geometric,fit,backgrounds}
\usepackage{titlesec}
\usepackage[most]{tcolorbox}
\usepackage{listings}
\usepackage{caption}
\usepackage{fancyhdr}
\usepackage{hyperref}

\definecolor{navy}{HTML}{10233F}
\definecolor{blue}{HTML}{2463A6}
\definecolor{teal}{HTML}{168C8C}
\definecolor{mint}{HTML}{E8F5F2}
\definecolor{pale}{HTML}{F2F6FA}
\definecolor{amber}{HTML}{A96308}
\definecolor{ink}{HTML}{233044}
\definecolor{graytext}{HTML}{58677B}
\definecolor{codebg}{HTML}{F3F5F7}

\hypersetup{
  colorlinks=true,linkcolor=blue,urlcolor=teal,
  pdftitle={Processus DevSecOps — Timesheet},
  pdfauthor={Safwen Amaira et Amine Mezlini},
  pdfsubject={Processus DevSecOps de l'IDE au déploiement}
}
\graphicspath{{images/}}
\setlength{\parskip}{0.4em}
\setlength{\parindent}{0pt}
\setlist[itemize]{leftmargin=1.35em,itemsep=0.2em,topsep=0.25em}
\setlist[enumerate]{leftmargin=1.55em,itemsep=0.25em,topsep=0.25em}
\renewcommand{\arraystretch}{1.2}
\captionsetup{font=small,labelfont={bf,color=navy},justification=centering}

\pagestyle{fancy}
\fancyhf{}
\fancyhead[L]{\small\color{graytext}Processus DevSecOps — Timesheet}
\fancyhead[R]{\small\color{graytext}ESPRIT • 5NIDS5}
\fancyfoot[L]{\small\color{graytext}Safwen Amaira • Amine Mezlini}
\fancyfoot[R]{\small\color{graytext}\thepage}
\renewcommand{\headrulewidth}{0.35pt}
\renewcommand{\footrulewidth}{0pt}

\titleformat{\chapter}[display]
  {\normalfont\bfseries\color{navy}}
  {\filleft\Large\chaptername\;\thechapter}
  {0.5ex}{\titlerule\vspace{1ex}\Huge}
  [\vspace{0.5ex}\titlerule]
\titleformat{\section}{\Large\bfseries\color{blue}}{\thesection}{0.65em}{}
\titleformat{\subsection}{\large\bfseries\color{teal}}{\thesubsection}{0.65em}{}

\newtcolorbox{infobox}[1][]{
  enhanced,breakable,colback=pale,colframe=blue!55!navy,
  boxrule=0.7pt,arc=2mm,left=2.5mm,right=2.5mm,top=1.5mm,bottom=1.5mm,#1
}
\newtcolorbox{resultbox}[1][]{
  enhanced,breakable,colback=mint,colframe=teal,
  boxrule=0.8pt,arc=2mm,left=2.5mm,right=2.5mm,top=1.5mm,bottom=1.5mm,#1
}
\newcommand{\pathfile}[1]{\nolinkurl{#1}}
\newcommand{\ProcessScreenshot}[3]{%
  \begin{figure}[H]\centering
  \IfFileExists{images/own/#1}{%
    \includegraphics[width=0.92\textwidth,height=7cm,keepaspectratio]{images/own/#1}%
  }{%
    \begin{tcolorbox}[width=0.92\textwidth,colback=pale,colframe=teal,
      boxrule=0.8pt,arc=2mm,height=4.1cm,valign=center]
      \centering\Large\color{navy}\textbf{Capture personnelle à insérer}\par\vspace{0.5em}
      \normalsize\color{ink}#2\par\vspace{0.6em}
      \small Fichier attendu : \pathfile{images/own/#1}
    \end{tcolorbox}%
  }%
  \caption{#3}\end{figure}%
}

\lstdefinestyle{terminal}{
  basicstyle=\ttfamily\small,backgroundcolor=\color{codebg},
  frame=single,rulecolor=\color{blue!30},breaklines=true,
  keepspaces=true,showstringspaces=false,columns=fullflexible,
  xleftmargin=0.4em,framexleftmargin=0.3em
}
\lstset{style=terminal}

\begin{document}

% Page de garde
\begin{titlepage}
\begin{tikzpicture}[remember picture,overlay]
  \fill[navy] (current page.north west) rectangle ([yshift=-8.7cm]current page.north east);
  \fill[teal] ([yshift=-8.7cm]current page.north west) rectangle ([yshift=-8.88cm]current page.north east);
  \fill[blue!12] ([yshift=1.5cm]current page.south west) rectangle (current page.south east);
  \draw[white,opacity=.12,line width=24pt] ([xshift=14.8cm,yshift=-1.4cm]current page.north west) circle (3.3cm);
  \draw[teal!70,opacity=.8,line width=5pt] ([xshift=14.8cm,yshift=-1.4cm]current page.north west) circle (2.25cm);
\end{tikzpicture}
\thispagestyle{empty}
\vspace*{1.1cm}
{\color{white}\Large\bfseries ESPRIT \hfill DÉPARTEMENT INFORMATIQUE}\par
\vspace{1.55cm}
{\color{teal!35!white}\large PROJET DE SÉCURITÉ DU PIPELINE CI/CD}\par
\vspace{0.35cm}
{\color{white}\Huge\bfseries DevSecOps}\par
\vspace{0.2cm}
{\color{white}\LARGE\bfseries Processus de sécurité continue}\par
\vspace{0.45cm}
{\color{white}\large De l'IDE au déploiement et à l'observabilité}\par

\vspace{2.15cm}
\begin{tcolorbox}[colback=white,colframe=white,arc=2mm,boxrule=0pt,width=0.89\textwidth,
  left=5mm,right=5mm,top=4mm,bottom=4mm]
\begin{tabularx}{\linewidth}{@{}>{\bfseries\color{navy}}lX@{}}
Réalisé par & Safwen Amaira et Amine Mezlini \\
Classe & 5NIDS5 \\
Encadrant & Taoufik Bessrour \\
Année universitaire & 2026–2027 \\
\end{tabularx}
\end{tcolorbox}
\vfill
{\color{navy}\large Rapport de processus et guide d'exploitation}\par
{\color{graytext}\small Projet Timesheet — DevSecOps-ESPRIT}\par
\end{titlepage}

\pagenumbering{roman}
\chapter*{Résumé}
\addcontentsline{toc}{chapter}{Résumé}
Ce rapport décrit uniquement le processus DevSecOps du projet Timesheet. Il suit le changement depuis l'analyse dans l'IDE, les contrôles pre-commit et l'intégration continue jusqu'à la validation Jenkins, au déploiement dans k3s, au test DAST et à la visualisation des métriques.

Le processus applique le principe \emph{shift-left} : les développeurs reçoivent des retours avant le push, tandis que les étapes CI/CD répètent les contrôles dans un environnement commun. Les règles bloquantes portent notamment sur les tests, les secrets, les erreurs SAST et les vulnérabilités critiques. Les rapports d'analyse restent consultables pour le triage et la correction.

\begin{infobox}
\textbf{Fil conducteur :} chaque étape vérifie son propre périmètre, conserve un résultat exploitable et ne transmet le changement à l'étape suivante qu'après satisfaction de ses conditions de passage.
\end{infobox}

\tableofcontents
\listoffigures
\clearpage
\pagenumbering{arabic}

% -----------------------------------------------------------------------------
\chapter{Le processus DevSecOps}

\section{Objectif et périmètre}
Le pipeline protège l'application Java/Spring Boot Timesheet et son déploiement. Les contrôles couvrent le code source, les dépendances Maven, les secrets, l'image conteneur, les fichiers d'infrastructure, l'application démarrée et l'environnement Kubernetes.

Le flux est automatisé à deux niveaux : les hooks du dépôt fournissent un premier filtre local au moment du commit; GitHub Actions contrôle les pushes et pull requests; Jenkins ne déploie qu'un commit dont le workflow GitHub a réussi. L'application déployée est ensuite vérifiée et sondée par OWASP ZAP.

\section{Vue complète du flux}
\begin{figure}[H]
\centering
\begin{tikzpicture}[
  node distance=4.5mm,
  stage/.style={draw=blue!55!navy,rounded corners=2mm,fill=pale,
    text width=0.80\textwidth,align=center,minimum height=8mm,
    inner sep=2.2mm,font=\small},
  gate/.style={draw=teal,rounded corners=2mm,fill=mint,
    text width=0.80\textwidth,align=center,minimum height=8mm,
    inner sep=2.2mm,font=\small},
  arr/.style={-{Stealth[length=2.2mm]},line width=0.8pt,color=teal!80!black}
]
\node[stage] (ide) {\textbf{IDE}\; SonarQube for IDE : analyse locale pendant l'édition};
\node[stage,below=of ide] (hook) {\textbf{Commit}\; pre-commit : hygiène, Gitleaks, Semgrep Java, Trivy IaC/SCA};
\node[gate,below=of hook] (gha) {\textbf{GitHub Actions}\; tests Maven, Gitleaks historique, Semgrep, Trivy, build image et ZAP};
\node[gate,below=of gha] (jenkins) {\textbf{Gate Jenkins}\; vérifie que le même SHA a un workflow GitHub réussi};
\node[stage,below=of jenkins] (security) {\textbf{Contrôles avant déploiement}\; scans secrets, SAST, dépendances, IaC, image et packaging};
\node[stage,below=of security] (deploy) {\textbf{Déploiement k3s}\; image, MySQL, Timesheet, rollout et healthcheck};
\node[stage,below=of deploy] (observe) {\textbf{DAST et suivi}\; OWASP ZAP, Pushgateway, Prometheus, Grafana};
\draw[arr] (ide) -- (hook);
\draw[arr] (hook) -- (gha);
\draw[arr] (gha) -- (jenkins);
\draw[arr] (jenkins) -- (security);
\draw[arr] (security) -- (deploy);
\draw[arr] (deploy) -- (observe);
\end{tikzpicture}
\caption{Étapes et quality gates du processus DevSecOps Timesheet.}
\label{fig:process}
\end{figure}

% -----------------------------------------------------------------------------
\chapter{Contrôles par étape}

\section{1 — IDE : prévenir pendant le développement}
SonarQube for IDE (anciennement SonarLint) donne un feedback de qualité et de sécurité pendant l'édition. Le mode connecté relie l'IDE au serveur SonarQube du projet; il requiert une liaison au projet et un jeton conservé dans le gestionnaire de secrets de l'IDE. Le développeur corrige les problèmes avant de créer le commit.

\ProcessScreenshot{01-ide-feedback.png}{Afficher l'analyse SonarQube for IDE sur un fichier Java du projet, avec le nom du projet visible.}{Feedback de sécurité dans l'IDE — capture à fournir par l'équipe.}

\section{2 — Pre-commit : filtrer le commit local}
Le framework \texttt{pre-commit} exécute les contrôles sur les changements staged. Les hooks vérifient l'hygiène des fichiers, les clés privées, les secrets avec Gitleaks, des règles Java Semgrep et les résultats Trivy IaC/SCA selon les fichiers touchés. Un échec bloque la création du commit jusqu'à correction ou exemption approuvée.

\begin{table}[H]
\centering
\caption{Contrôles locaux du commit}
\begin{tabularx}{\textwidth}{@{}>{\bfseries}p{3.1cm}p{3.3cm}X@{}}
\toprule
Contrôle & Outil & Déclencheur et résultat attendu \\
\midrule
Hygiène & pre-commit-hooks & À chaque commit; corrige ou bloque les erreurs de format, fichiers volumineux et clés privées. \\
Secrets staged & Gitleaks & Analyse le contenu staged; une détection empêche le commit. \\
SAST Java & Semgrep & Analyse les fichiers Java concernés; les règles de sécurité en erreur bloquent le commit. \\
IaC & Trivy config & Lorsqu'un fichier IaC ciblé change; bloque selon le seuil configuré dans le hook. \\
Dépendances & Trivy fs & Lorsque le manifeste Maven change; bloque les vulnérabilités critiques selon les options du hook. \\
\bottomrule
\end{tabularx}
\end{table}

\ProcessScreenshot{02-precommit-pass.png}{Montrer un commit déclenchant les hooks avec leurs statuts; si possible, ajouter une capture d'un finding bloquant corrigé.}{Résultat pre-commit — capture personnelle à fournir.}

\section{3 — GitHub Actions : contrôles à chaque changement distant}
Le workflow \texttt{.github/workflows/devsecops.yml} se déclenche sur \texttt{push}, \texttt{pull\_request} et \texttt{workflow\_dispatch}. Il prépare Java 17 et MySQL, exécute Maven, analyse l'historique complet avec Gitleaks, lance Semgrep, Trivy SCA/IaC/image, construit l'image, démarre l'application et exécute un baseline ZAP. Les rapports sont uploadés sous forme d'artefact pour le run.

Les seuils de blocage sont explicites : tests Maven en échec, findings Gitleaks, findings Semgrep de sévérité ERROR et vulnérabilités CRITICAL non ignorées bloquent le workflow. Les rapports Trivy incluent aussi les sévérités HIGH; la politique CI utilise le seuil CRITICAL pour le blocage. ZAP s'exécute lorsque les tests et le build d'image ont réussi; les options du baseline déterminent le traitement des warnings.

\ProcessScreenshot{03-github-actions.png}{Capturer le résumé du workflow GitHub Actions et ses étapes SAST, SCA, IaC, image, DAST et upload d'artefact.}{GitHub Actions — capture personnelle d'un run du dépôt.}

\section{4 — Jenkins : exiger le résultat du même commit}
Le job Jenkins est provisionné automatiquement et interroge GitHub toutes les cinq minutes. Il récupère \texttt{master}, calcule son SHA et recherche un run \texttt{push} du workflow \texttt{DevSecOps} portant exactement ce SHA. En l'absence d'un résultat terminé avec succès, Jenkins marque le build non exécuté et ne déploie pas.

Après le gate, Jenkins répète Gitleaks, Semgrep, Maven, Trivy dépendances/IaC et le scan d'image avec des seuils critiques. Les rapports sont archivés. Cette répétition protège contre une divergence entre les contrôles CI et l'étape de déploiement.

\ProcessScreenshot{04-jenkins-gate.png}{Montrer la page Jenkins du job provisionné, le SHA candidat et les stages réellement exécutés.}{Gate Jenkins et stages — capture personnelle d'un build récent.}

\section{5 — Déploiement contrôlé dans k3s}
Après les gates, Jenkins applique les manifestes Kubernetes de la base et de l'application, crée le secret de base à partir du credential d'environnement, importe l'image dans le runtime containerd de k3s, met à jour l'image du déploiement et attend les rollouts. Un contrôle HTTP confirme ensuite la santé de l'application.

Le cluster est lancé localement par Docker Compose. Sur un hôte où la commande \texttt{docker} vise Podman rootless, le projet fournit \pathfile{scripts/compose-rootful.sh} afin d'utiliser le socket Docker rootful, nécessaire à k3s dans cet environnement.

\ProcessScreenshot{05-kubernetes-rollout.png}{Capturer les pods Timesheet et MySQL prêts, ainsi que la sortie du rollout ou du healthcheck.}{Déploiement Kubernetes — capture personnelle d'un run récent.}

\section{6 — DAST puis métriques}
OWASP ZAP effectue un scan baseline sur l'application déployée. Jenkins conserve les exports et publie les informations de réussite/échec vers Pushgateway. Prometheus collecte et historise les métriques; Grafana affiche le dashboard provisionné du projet.

\ProcessScreenshot{06-zap-report.png}{Capturer le résumé du rapport ZAP du run présenté, avec la cible et la date visibles; masquer toute donnée sensible.}{DAST OWASP ZAP — capture personnelle.}

\ProcessScreenshot{07-grafana-dashboard.png}{Capturer le dashboard Grafana DevSecOps montrant l'historique des builds et les résultats DAST.}{Suivi des résultats dans Grafana — capture personnelle.}

% -----------------------------------------------------------------------------
\chapter{Gates, rapports et exploitation}

\section{Politique de progression}
\begin{table}[H]
\centering
\caption{Conditions principales de passage}
\begin{tabularx}{\textwidth}{@{}>{\bfseries}p{3.0cm}p{5.0cm}X@{}}
\toprule
Étape & Condition de passage & En cas d'échec \\
\midrule
Commit local & Hooks applicables terminés sans erreur bloquante & Corriger le fichier ou le finding; le scan qui ne démarre pas doit être réparé et relancé. \\
GitHub Actions & Tests, secrets, SAST et gates Trivy sous leurs seuils & Le workflow échoue; consulter les logs et les artefacts du run. \\
Jenkins & Workflow GitHub réussi pour le SHA exact de master & Aucun déploiement; attendre les checks ou corriger et repousser. \\
Déploiement & Rollouts Kubernetes et endpoint de santé réussis & Le build échoue avant la phase DAST; examiner les événements et logs Kubernetes. \\
DAST & Baseline ZAP terminé selon les seuils configurés & Examiner les alertes, corriger ou appliquer une exemption tracée. \\
\bottomrule
\end{tabularx}
\end{table}

\section{Rapports et notifications}
Les contrôles produisent des rapports JSON, et ZAP produit un rapport HTML/JSON. GitHub Actions attache les rapports à l'artefact \texttt{devsecops-reports-<run id>}; Jenkins archive \texttt{reports/**}. Les métriques suivent les builds dans Prometheus/Grafana. Une notification Slack peut être activée avec un webhook fourni hors dépôt via la configuration locale.

Les secrets locaux sont conservés dans un fichier \texttt{.env} ignoré par Git. Dans un environnement partagé, utiliser des credentials protégés de Jenkins/GitHub ou un gestionnaire de secrets. Ne jamais inscrire une clé ou un mot de passe dans le Jenkinsfile, les journaux, les captures ou les rapports publiés.

\section{Procédure d'exemption}
Une exemption ne s'applique qu'à un faux positif confirmé ou à un risque accepté temporairement. Elle doit référencer une issue et préciser l'outil/règle ou CVE, le fichier/composant, les preuves, l'impact, les mesures compensatoires, le responsable, l'approbateur et une date d'expiration. La suppression est la plus étroite possible et est revue par une personne différente de l'auteur. Un secret potentiellement exposé doit être révoqué et renouvelé; une exemption ne remplace pas sa rotation.

% -----------------------------------------------------------------------------
\chapter{Utilisation reproductible}

\section{Préparer le poste et démarrer}
Prérequis : Git, Java 17, Maven, pre-commit, IntelliJ avec SonarQube for IDE et Docker Engine rootful avec Compose. Configurer le mode connecté Sonar si le serveur local est démarré; placer le jeton dans le gestionnaire sécurisé de l'IDE.

\begin{lstlisting}[language=bash,caption={Installer les hooks et démarrer le laboratoire}]
./scripts/setup-devsecops.sh
mvn --batch-mode test

# Start the stack with the Docker engine selected by the project
./scripts/compose-rootful.sh up -d --build
./scripts/compose-rootful.sh ps -a
\end{lstlisting}

\section{Suivre une modification}
\begin{enumerate}
  \item Examiner les retours IDE et corriger le code.
  \item Lancer tests et scanners ciblés, puis créer le commit avec les hooks actifs.
  \item Pousser la branche ou ouvrir une pull request; suivre le workflow GitHub.
  \item Après le succès sur le commit de \texttt{master}, suivre le prochain cycle Jenkins (polling toutes les cinq minutes).
  \item Examiner le rollout k3s, le résultat ZAP et le dashboard Grafana.
  \item Archiver les rapports et noter le SHA, la date et les exceptions approuvées.
\end{enumerate}

\section{Services locaux}
\begin{table}[H]
\centering
\caption{Interfaces de l'environnement de démonstration}
\begin{tabularx}{\textwidth}{@{}>{\bfseries}p{3.6cm}p{5.1cm}X@{}}
\toprule
Service & Adresse locale & Usage dans le processus \\
\midrule
Jenkins & \url{http://localhost:8080} & Gate et livraison \\
SonarQube & \url{http://localhost:9000} & Analyse IDE connectée \\
Timesheet & \url{http://localhost:8082/timesheet-devops/} & Application locale et DAST Compose \\
Prometheus & \url{http://localhost:9090} & Historique des métriques \\
Grafana & \url{http://localhost:3000} & Tableau de bord \\
k3s & \url{https://localhost:6443} & Cluster local \\
\bottomrule
\end{tabularx}
\end{table}

% -----------------------------------------------------------------------------
\appendix
\chapter{Remplacer les captures par vos propres images}
Le rapport ne contient aucune capture issue du dossier screenshots du dépôt. Les cadres visibles dans le PDF sont des emplacements à remplir avec des captures prises par l'équipe pendant ses propres exécutions. Le schéma de pipeline est dessiné directement en TikZ dans ce document.

Pour ajouter une image, copier le fichier dans \pathfile{images/own/} en utilisant exactement l'un des noms ci-dessous. Le rapport détecte le fichier et remplace automatiquement le cadre par la capture. Les images doivent montrer le projet/résultat pertinent et masquer les mots de passe, tokens, adresses privées et données personnelles.

\begin{longtable}{@{}p{4.2cm}p{5.1cm}p{5.0cm}@{}}
\caption{Noms attendus et contenu conseillé pour les captures personnelles}\label{tab:screenshots}\\
\toprule
Nom du fichier & Capture à prendre & Figure \\
\midrule
\endfirsthead
\toprule Nom du fichier & Capture à prendre & Figure \\
\midrule\endhead
\bottomrule\endfoot
\pathfile{01-ide-feedback.png} & IntelliJ avec une analyse SonarQube for IDE sur le code Timesheet. & Figure~\ref{fig:process-ide} \\
\pathfile{02-precommit-pass.png} & Terminal montrant les hooks pre-commit exécutés sur votre changement. & Figure~\ref{fig:process-hook} \\
\pathfile{03-github-actions.png} & Run GitHub Actions et ses stages de sécurité. & Figure~\ref{fig:process-gha} \\
\pathfile{04-jenkins-gate.png} & Jenkins avec le SHA et les stages d'un build récent. & Figure~\ref{fig:process-jenkins} \\
\pathfile{05-kubernetes-rollout.png} & Pods prêts et rollout du déploiement Timesheet. & Figure~\ref{fig:process-k8s} \\
\pathfile{06-zap-report.png} & Rapport DAST ZAP du même déploiement. & Figure~\ref{fig:process-zap} \\
\pathfile{07-grafana-dashboard.png} & Dashboard Grafana avec l'historique des builds/scans. & Figure~\ref{fig:process-grafana} \\
\bottomrule
\end{longtable}

\section*{Compiler dans Overleaf}
\addcontentsline{toc}{section}{Compiler dans Overleaf}
Créer un projet Overleaf depuis le ZIP fourni ou téléverser \texttt{rapport-devsecops.tex} comme document principal avec le dossier \texttt{images/own/}. Choisir pdfLaTeX et compiler deux fois pour actualiser les références. Le ZIP contient \texttt{main.tex} à sa racine et les emplacements d'images; ajouter ensuite vos fichiers dans \texttt{images/own/}.

\section*{Structure des fichiers du processus}
\addcontentsline{toc}{section}{Structure des fichiers du processus}
\begin{longtable}{@{}p{5.3cm}p{8.7cm}@{}}
\toprule
Chemin & Rôle dans le processus DevSecOps \\
\midrule
\pathfile{.pre-commit-config.yaml} & Définition des contrôles de commit local \\
\pathfile{scripts/precommit-*.sh} & Lancement conteneurisé des scanners locaux \\
\pathfile{.github/workflows/devsecops.yml} & Pipeline CI sur push et pull request \\
\pathfile{pipeline/Jenkinsfile} & Gate du SHA GitHub, contrôles, déploiement et DAST \\
\pathfile{docker-compose.yml} & Services locaux du processus \\
\pathfile{kubernetes/} & Ressources de déploiement Timesheet/MySQL \\
\pathfile{monitoring/} & Provisioning Prometheus/Grafana \\
\pathfile{security/semgrep-rules.yml} & Règles SAST propres au projet \\
\pathfile{reports/} & Rapports JSON/HTML générés par les scans \\
\bottomrule
\end{longtable}

\vfill
\begin{center}
\color{graytext}\small Safwen Amaira • Amine Mezlini • ESPRIT 5NIDS5
\end{center}
\end{document}
